\begin{mistake}{2026-06-15}{02}

\begin{problemcontent}
设函数 \(f(x)\) 有二阶连续导数，且 \(f(0)=0\)，\(f'(0)>0\)，\(f''(0)<0\)，则
(A) \(x=0\) 是 \(|f(x)|\) 的极值点，但 \((0,f(0))\) 不是曲线 \(y=|f(x)|\) 的拐点.
(B) \(x=0\) 不是 \(|f(x)|\) 的极值点，但 \((0,f(0))\) 是曲线 \(y=|f(x)|\) 的拐点.
(C) \(x=0\) 是 \(|f(x)|\) 的极值点，且 \((0,f(0))\) 是曲线 \(y=|f(x)|\) 的拐点.
(D) \(x=0\) 不是 \(|f(x)|\) 的极值点，且 \((0,f(0))\) 不是曲线 \(y=|f(x)|\) 的拐点.
\end{problemcontent}

\begin{solutioncontent}
正确答案：(C)

令 \(g(x) = |f(x)|\)。
1. \textbf{极值判定}：已知 \(f(0)=0\)，故 \(g(0) = |f(0)| = 0\)。由于绝对值非负，对任意 \(x\) 恒有 \(g(x) \ge 0\)，即 \(g(x) \ge g(0)\)。由极值定义，\(x=0\) 是 \(g(x)\) 的极小值点。
2. \textbf{拐点判定}：已知 \(f(0)=0\) 且 \(f'(0)>0\)，由局部保号性，在 \(x=0\) 邻域内：
当 \(x<0\) 时，\(f(x)<0\)，\(g(x)=-f(x)\)；当 \(x>0\) 时，\(f(x)>0\)，\(g(x)=f(x)\)。
求二阶导 (\(x \neq 0\))：
\[
g''(x) = \begin{cases} 
-f''(x), & x < 0 \\ 
f''(x), & x > 0 
\end{cases}
\]
由于 \(f''(0)<0\) 且具有连续性，在 \(x=0\) 邻域内恒有 \(f''(x)<0\)。
因此当 \(x<0\) 时，\(g''(x)=-f''(x)>0\)（曲线下凹）；当 \(x>0\) 时，\(g''(x)=f''(x)<0\)（曲线上凸）。
因为 \(g(x)\) 在 \(x=0\) 处连续，且两侧二阶导数变号（凹凸性改变），故 \((0,0)\) 是曲线 \(y=|f(x)|\) 的拐点。
综上，选 (C)。
\end{solutioncontent}

\mistakeknowledge{
\begin{itemize}
  \item \textbf{绝对值极小值}：遇到 \(g(x)=|f(x)|\) 且 \(f(x_0)=0\) 时，利用非负性 \(|f(x)| \ge 0\) 即可直接判断 \(x_0\) 为极小值点，无需借助导数。
  \item \textbf{易错点（不可导点作拐点）}：考研大纲中，拐点仅要求是“连续”曲线的凹凸分界点，不要求在该点可导。因图像翻折产生的尖点也可以是拐点。
  \item \textbf{图像翻折的凹凸性}：函数图像在 \(x\) 轴下方部分被加上绝对值翻折后，其凹凸性必然发生反转（如上凸变下凹）。
\end{itemize}}

\end{mistake}
